% 原书第 99 页（PDF 第 112 页）
\BookPage{99}{112}
\BookHeading{清乾隆帝南巡与扬州}{清乾隆帝南巡与扬州}{岑大利*}
扬州是长江下游南北运河沿线的一个重要城市，有着上千年的悠久历史，风景秀丽，园林众多，既是工商业发达的地区，也是文人荟萃的风雅之地。清朝时，乾隆帝六次南巡都经过这里，留下了许多宝贵的历史遗迹。

\BookSubheading{一、乾隆南巡至扬州的路线}

乾隆帝南巡的时间和路线与康熙帝大致相同：每次南巡都于正月十五日前后从京师出发，陆路经直隶、山东到江苏的清口渡黄河，乘船沿运河南下，经扬州、镇江、丹阳、常州、苏州进入浙江境内，再由嘉兴、石门抵达杭州。回銮时，绕道江宁，祭奠明太祖陵，于四月下旬或五月初返回京师。这条路线告诉我们，扬州是乾隆帝在进入江南之后，走水路，沿运河南下所经过的城市。

乾隆帝到扬州后，驻跸多日，据记载，其巡幸的路线有三站，第一站自北桥始，走香阜寺（乾隆赐名）御道，到天宁寺行宫，计62里；第二站自天宁寺行宫入天宁门，出钞关码头登船，经文峰寺，九龙桥，到高旻寺行宫，计16里；第三站由高旻寺行宫走锦春园，经陈家湾，经由闸，到江口，计23里。这是在天宁寺行宫修成之后的巡幸路线。在天宁寺行宫修成之前，也就是十六年和二十二年的南巡，是由香阜寺入天宁门出钞关码头，然后至塔湾，到蜀冈三峰及黄、江、程、洪、张、汪、周、王、闵、吴、徐、鲍、田、郑、巴、余、罗、尉诸园亭，或便道，或于塔湾迂道。乾隆走水路乘御舟，御舟前有两对船走两边引路，船旁令一人骑马在河边行走，以备差遣。拉绎的用州县盐块，不够则雇民夫。另外有四只大船名安福舻、

\vfill
\begin{BookNotes}
\BookNote{*}{作者任职于中共中央党校文史部历史教研室。}
\BookNote{1.}{李斗《扬州画舫录·草河录上》卷1，第1-3页，中华书局，2001年版。}
\end{BookNotes}

% 原书第 100 页（PDF 第 113 页）
\BookPage{100}{113}
\noindent 翔凤艇、湖船、扑拉船，跟在其后，其余船则装载各种用品及随从官兵。船行时，河的两岸及桥头村口，各安排士兵守卫，禁止私船出入，绎道每里安设兵丁三名，令村镇民妇跪伏瞻仰，御舟经过需回避时，令男子退出，不禁妇女。2

乾隆在扬州巡幸时所走的御道，是用文砖（或用石板）铺成的，帮宽三尺，中心正路一丈六尺，两傍马里各七尺，都要求坚实平整笔直，不能随意弯曲，为此，许多民宅被拆毁，坟墓被挖掘，良田被毁坏。此外，所经之路都要用黄土铺垫，泼水清尘。道旁或搭彩棚、或陈水嬉，并动员老百姓沿路迎接。乾隆三十年，游上方寺时，由于「万民随马足趋瞻」，出现了「或有践踏麦苗者」的景象，于是乾隆作诗云：「马足纷随定何碍，蹂躏惟惜麦苗芒。」3

\BookSubheading{二、扬州盐商对乾隆南巡的贡献}

乾隆南巡的财政支柱主要是两淮的盐商，而两淮盐运使司公署就设在扬州城内，八大商总也始终住在扬州城内。两淮商人究竟掌握了多少资本，具体数字不得而知，但据史料记载，「扬州盐商豪侈甲天下，百万以下者皆谓之小商」。这些商人尽管经济上富逾王侯，但在政治上却无权无势，他们对封建官府存在着严重的依附关系。因此，每次南巡时，当地官员都要从他们手中勒索一部分资金，作为自己巴结、讨好皇帝之用，更重要的还是让他们筹办各种南巡工程，如营建行宫、园林等，而这些富商也正想借南巡讨好官府，特别是皇帝，为自己赢得荣誉和权势。乾隆也深知这一点，对他们的这种做法，不但没有加以制止，反而给予赞扬，他说：「今两淮商人踊跃急公，捐输报效，地方官一应公务，俱于此取给，可信其无丝毫扰累闾阎，而办理工程，贫民更得其资其利，赖于生计，转有裨益。」4这更加鼓舞了商人们的贡献积极性。

盐商对乾隆南巡的资助首先是捐献银子。如二十六年、四十五年、四十九年，两淮盐商每次都捐银100万两，这些钱中的相当一部分被充作南巡时扬州官府办差

\vfill
\begin{BookNotes}
\BookNote{2.}{《清稗类钞·巡幸类》第1册，第341页，「高宗南巡禁卫之严」条，中华书局，1984年版。}
\BookNote{3.}{李斗，《扬州画舫录·草河录上》卷1，第3页。}
\BookNote{4.}{清高晋等纂辑，《南巡盛典·褒赏》卷6，第6页。乾隆三十六年刻本。}
\end{BookNotes}

% 原书第 101 页（PDF 第 114 页）
\BookPage{101}{114}
\noindent 及乾隆赏赐之用。此外，他们还进献礼物，有诗、赋、字画和文集，也有进呈和时事有关的论文。乾隆对此都赏赐从优。

除了捐银、进献礼物之外，盐商还直接办理接驾等事宜。如南巡时，盐商负责供应皇帝的日常所需，因巡幸队伍异常庞大，这笔开支为数不小。如准备薪炭的商人，地方官命令他们加意供给，薪炭虽堆积如山，但顷刻就使用无余；准备清理粪便的商人，地方官命令他们沿塘遍置盆盂，上加木盖，以备绎夫们的便溺，每县动辄就千万个，而御舟一过，马上就被役夫们破坏，很少有完好的，等御舟再经过，又需要置备新的盆盂。

盐商还出资修建行宫，据史载，行宫在扬州有四，一在金山，一在焦山，一在天宁寺，一在高旻寺。其中，高旻寺行宫是康熙南巡时修建的，而天宁寺行宫则是专为乾隆南巡而修建的。该寺居扬州八大寺（天宁、建隆、重宁、慧因、法浄、高旻、静慧、福缘）之首，乾隆二十一年由商人谷裕丰、张世天、尉宏美等出资在寺的右边建起行宫。据《南巡盛典》所载的「天宁寺行宫图」所示，该行宫有宫门、大宫门、二宫门、前殿、寝殿；右宫门，戏台、前殿。垂花门、寝殿、西殿、内殿、御花园。门前左右茶膳房，两边为护卫房。最后为后门，通重宁寺。可以说是皇帝起居、办公、娱乐的设施一应俱全。乾隆驻跸过后，该行宫的各门都档上木栅，禁止游人入内。在乾隆南巡期间，扬州所建行宫的出资人数量，据《扬州画舫录》记载，有各种商人共十八姓，其中多数是盐商。而据第一历史档案馆所藏《扬州行宫名胜全图》所标明的集资修建行宫的商人在二十五姓以上。如香阜寺是商人王履泰出资，建楼廊209间，亭台17座。这二十五姓一共修建了楼廊5154间，亭台196座。

行宫内的布置和陈设都非常奢华，各种书籍、字画、瓷器、香炉、挂屏等玩好器皿，应有尽有，不计其数。这些玩好器物大多是地方督抚向通省大小官员临时借用或重金购买的，乾隆对此是心知肚明，并担心由此而造成种种情弊，所以在二十六年下令，「除高旻寺行宫原系商人自行置办，应仍听其预备，但须交收清楚

\vfill
\begin{BookNotes}
\BookNote{5.}{李斗，《扬州画舫录·新城北录中》卷4，第102页。}
\BookNote{6.}{第一历史档案馆，《扬州行宫名声全图》。}
\end{BookNotes}

% 原书第 102 页（PDF 第 115 页）
\BookPage{102}{115}
\noindent 外，其余各处行宫，概无得陈列玩器」7。但实际上，接驾的官员为讨好皇帝，仍然我行我素。行宫院落中还需要数量巨大的花草竹木、珍石和水池、细桥等种种设施。如乾隆第一次南巡时，扬州平山堂行宫本来没有梅花，盐商们觉得风景不佳，竟捐资植梅一万株，以作点缀之用。

盐商们还广建园林，逢迎乾隆，「如汪氏之净香园，黄氏之趣园，洪氏之倚虹园，汪氏之九峰园等」，「当时各大商所造园林，甲于天下」8。据说两淮八大盐商之一汪石公死后，家务由其妻子主持，人称汪太太。乾隆将巡幸扬州，汪太太和其他盐商提前在城郊择荒地数百亩，仿照杭州西湖风景，建筑亭台楼阁，可就是缺少一个湖泊，眼看巡幸将至，汪太太毅然独出数万金，招集工匠，修造「三仙池」。水池完工的第二天，乾隆就到了，他游览后非常满意，对汪太太大加赞扬，赏赐珍物。于是这位商人之妇一时名闻远近。更有甚者，一夜建造一座白塔。《清稗类钞》记述道：「高宗巡幸至扬州，时江某为盐商总，承办一切供应。某日，高宗幸大虹园，至一处，顾左右，曰:此处颇似南海之琼岛春阴，惜无塔耳。江闻之，即以万金贿近侍，图塔状，既得图，乃鸠工庀材，一夜而成。次日又幸园，建塔巍然，大异之。以为伪也，即之，果砖石所成，询之其故，叹曰：盐商之财力伟哉！乾隆六次南巡，扬州盐商一次又一次地建造园林，园林加上行宫，从扬州城里到平山堂，处处亭园楼台，号称「一路楼台直到山」。乾隆面对这些曲水回栏，崇楼杰阁，也不觉赞叹盐商们资力雄厚。

盐商们为迎接乾隆南巡，还不惜花费巨额财富，蓄养花、雅两部戏曲班子，从高桥起，到迎恩桥上，在数十里长的河两岸，搭起戏台，奏乐演戏。如在天宁寺前搭起松棚做戏台，演仙佛麟凤太平鼓等戏，叫做大戏，事情完后就拆卸。后来重宁寺搭建了大戏台，就把演大戏搬到了那里。1据说曾有新剧《雷锋塔传奇》，因时间紧，来不及排练，于是让演员用旧调唱新词。演出时，演员偶尔忘了台词，就含

\vfill
\begin{BookNotes}
\BookNote{7.}{《南巡盛典·恩伦》卷2，第5页。}
\BookNote{8.}{王振世，《扬州揽胜录·湖山识略》。}
\BookNote{9.}{《清稗类钞·园林类》第1册，第206页，「大虹园之塔」条。}
\BookNote{10.}{李斗，《扬州画舫录·新城北录下》卷5，第107页。}
\end{BookNotes}

% 原书第 103 页（PDF 第 116 页）
\BookPage{103}{116}
\noindent 混唱过。当乾隆乘船航行时，演员就在戏台上面对御舟演唱，这种别出心裁的戏剧表演往往使乾隆感到新奇、高兴，但乾隆也担心如此铺张，会带来不好的影响，所以在乾隆二十年下令「沿途水次从前俱设灯船、戏船、台阁俚俗游玩之具，此不过地方吏役及民间游手好事之徒藉名苛敛，尤当通行严禁」12。但因地方官讨好皇帝的心不变，所以这种演戏还是一直进行下去，由此也带动了扬州戏剧的发展。乾隆四十二年，巡盐御史伊龄阿奉旨在扬州设局修改曲剧，历经四年完成，留下了宝贵的遗产。

\BookSubheading{三、乾隆帝对扬州盐商的褒奖}

南巡中，乾隆的特殊褒奖种类繁多，题字最为常见。在扬州驻跸时，他曾书写许多福字分赐给官商缙绅。当时扬州正盛行修建私家园林之风，得到御书福字的人一个个喜气洋洋，纷纷盖起福字厅以相炫耀。乾隆南巡时驻跸的天宁寺，曾得到乾隆帝赐给的匾额七个，即「淮南香界」「浮山华海」「淮南丽曙」「神威拥护」「省方设教」「大雄宝殿」「万佛楼」。盈联八副，即「花雨南天，灵文传妙谛；香空蜀阜，藩墅表名区」；「闾里讴歌闻乐恺，轩窗烟景遍清嘉。」；「楚尾吴头开画镜，林光鸟语入吟轩。」；「定地生欢喜，香台普吉祥。」；「众香馥郁凝华盖，多宝光明驻法轮。」；「琉璃瓶水资功德，缨络龛云现吉祥。」；「西竺驻祥轮，三摹合证，东山留净业，二谛俱融。」「商鼎周彝自典重，槛花苑树相芬芳」。都供奉在大殿里，另有乾隆赐给的御制七言律诗四首，刻石供奉在碑亭中。13

乾隆有时也给商人一些实际好处，如四十五年，免除两淮盐商未交税银一百二十万两，缓征银二十七万两；四十七年，免除淮南商人未完银二百万两。此外，还允许盐商在规定领取的盐引之外，每引加耗二十斤。14通过这些应征课额的减免和加耗，乾隆博得「圣恩浩荡」的美名。但事实上，他的这些减免和加耗在数

\vfill
\begin{BookNotes}
\BookNote{11.}{《清稗类钞·巡幸类》第1册，第341页，「高宗南巡供应之盛」条。}
\BookNote{12.}{《南巡盛典·恩伦》卷2，第9页。}
\BookNote{13.}{李斗，《扬州画舫录·新城北录中》卷4，第86页。}
\BookNote{14.}{《南巡盛典·恩伦》卷1，第9页。《扬州府志·巡幸》卷1同。}
\end{BookNotes}

% 原书第 104 页（PDF 第 117 页）
\BookPage{104}{117}
\noindent 量上是非常有限的，而商人在巡幸期间的巨额开支，以及因政治上、经济上对封建官府严重依赖而造成的巨大损失，则是无法用数字进行统计的。

乾隆还授予他们大量的空头官衔。例如十六年南巡，乾隆发布上谕说，两淮商人「踊跃急公，捐输报效」，令将其职衔加顶戴一级；二十二年又以两淮众商，虽不是当官的人，但承办差务积极，令将官秩在三品以上者，赏给奉宸苑卿衔，未到三品者，各加顶戴一级。15二十七年二月又下谕旨，给予「已加奉宸苑卿衔之黄履暹、洪征治、江春、吴禧祖各加一级，已加按察使衔之徐士业、汪立德、王助俱着加奉宸苑卿衔，李志薰、汪秉德、毕本恕、汪寿着各加按察使衔，程征启着赏给六品职衔，程扬宗、程玲、吴山玉、汪长馨俱着各加一级」16。这些空头官衔没有什么实际意义，但可以荣耀其身分，而这正是商人们朝思暮盼、求之不得的。因此，一旦得到封典，他们便感恩戴德，「颂声如雷」。

四十九年南巡时，乾隆将朝廷缮写好的《四库全书》藏于扬州大观堂的文汇阁，并令地方大吏允许士子借阅。

\BookSubheading{四、南巡对扬州城市发展的影响}

乾隆时期，国家财力雄厚，统治阶级中浮华、享乐之风日益加剧，尤其是乾隆南巡时，大量商人、绅士、文人墨客、优伶杂技聚集于皇帝驻跸的几个城市，极大地刺激了这些城市的商品消费；与此同时，官府、名流为了巴结、讨好皇帝，大肆修缮、装饰这些城市的园林、道路、民居、商店、作坊。一时间扬州等中心城市急剧繁华。早在乾隆南巡以前，扬州就是商人汇聚之地，商品经济比较发达，南北商人多向这里聚集，所以扬州城内就形成了一些包括南北货在内的商业街，如缎子街，大量的绸缎铺户在这里开业。南矢巷、北矢市、南皮市、北皮市，彩衣街等都聚集了许多南北货物铺户。还有以南北名酒为货的酒楼集中在虹桥附近，通州的雪酒、泰州的枯酒、陈老枯酒、高邮木瓜酒、五加皮酒、宝应乔家白酒、绍兴老酒等南北各地名酒，都可在这里买到。南北商人的聚集，使酒楼、茶肆、戏场、妓院也

\vfill
\begin{BookNotes}
\BookNote{15.}{《扬州府志·巡幸》卷1。嘉庆十五年刻本。}
\BookNote{16.}{《南巡盛典·褒赏》卷6，第8页。}
\end{BookNotes}

% 原书第 105 页（PDF 第 118 页）
\BookPage{105}{118}
\noindent 随之发展起来。孔尚任就写过这样几句诗：「东南繁华扬州起，水陆物力盛罗绮。朱桔黄橙香者掾，蔗仙糖狮如茨比。一客已开十丈筵，客客对列成肆市。」其盛况可见一斑。

然而，市镇建设却相对落后，整个城市尚未脱离乡村集镇气象。据史书记载，乾隆十几年，这里还很荒凉，河道宽不过两丈，远远望去象长绳一般，河岸亭台很少，黄土白沙，杂草虫鸣。自从十六年南巡以后，商人雇工兴修，增荣饰观，于是河道变宽，水势浩浩，迂回曲折，其风景之秀丽迷人，竟非文人画士所能描述。但是这种繁华是和当时社会生产力的实际程度相脱节的，它只是由于清代帝王的南巡促成的，繁华中的病态因素在以后的历史中产生了不良影响，尤其是当乾隆停止巡幸以后，扬州城就呈现出衰落迹象。从嘉庆初年扬州便开始由盛转衰，楼台倾毁，花木凋零，此后荒芜更加严重，当初修造园林的商人大多歇业贫散，画馆寒士也多清苦，吏仆佣贩皆不能糊口。加上江淮水患，下河的饥民由楚黔至滇城结队乞食。至于南巡期间修造的园林楼阁，此时竟也是十之七八成为了黄土。


% 原书第 106 页（PDF 第 119 页）
\BookPage{106}{119}
